% discussion.tex -- Results: discussion (draft 4 Oct 2026). To be revised once the hardware detach,
% the four-UAV rig flights and the rig-twin attach are in.

\subsection{Discussion}
\label{sec:discussion}
% 9 Oct (structure audit): the list of simulator artefacts moved here from the experimental method;
% the steady-error taxonomy points at Table tab:error_sources.

This section answers the questions in order, then discusses what the results mean beyond them: the
role of geometry, the steady errors of a re-planned cascade, the twin against the rig, and the
limitations.

\paragraph{RQ1: a centralised controller through join and leave.} In simulation the answer for the
detach is yes, with a condition. A detach by OCP resize redistributes the load in the first solve after the release,
and the payload settles level within a few seconds, provided the survivors' attachment points still
surround its centre of mass. The centralised planner's advantage is that it knows the
whole attachment geometry. It does not have to discover the new load distribution from the motion
that follows the event: it solves for it. This supports the hypothesis of the project plan. Neither
reconfiguration has been flown on the rig.
% CHECK: stale since the 7-8 Oct rig detach flights (Section 4.3, fig:rig_detach).

\redo{the attach half of RQ1: the weld transient of an attach by OCP resize in the full rejoin mission (Section~\ref{sec:sim_attach}).}
% --- original text, flown before the rig twin; kept to refill this place:
% \paragraph{RQ1: a centralised controller through join and leave.} In simulation the answer is yes,
% with a condition. A detach by OCP resize redistributes the load in the first solve after the release,
% and the payload settles level within a few seconds, provided the survivors' attachment points still
% surround its centre of mass. An attach by OCP resize folds the newcomer in with a weld transient of
% under 2$^\circ$ in the full rejoin mission. The centralised planner's advantage is that it knows the
% whole attachment geometry. It does not have to discover the new load distribution from the motion
% that follows the event: it solves for it. This supports the hypothesis of the project plan. Two parts
% of the answer are not yet in. Neither reconfiguration has been flown on the rig, and the attach
% results come from the earlier simulation geometry, before the rig twin.

\redo{the attach's hand-over faults (mux authority on an armed but idle stream, the newcomer's starting tension, the frozen target, the weld gate's clock), restated once the attach has flown on the twin.}
% --- original text, flown before the rig twin; kept to refill this place:
% The attach needed much more than the resize. Most of the faults that blocked it were in the hand-over
% between two controllers, not in either controller. A multiplexer gave authority to a stream that was
% armed but idle. A newcomer started with too much tension pulled the ring with a lever about the
% incumbents' apex. A frozen target let the payload coast away from the approach. A weld gate measured
% speed in the wrong clock. A joining UAV is a hand-over problem as much as a control problem, and
% each interface needed an explicit rule for who commands the UAV and when.

\paragraph{RQ2: trajectory-based against dissipative.}
\redo{the answer to RQ2 (trajectory-based planner against the dissipative network in tracking, detach and attach), rewritten from the re-flown comparison of Section~\ref{sec:sim_architecture}.}
% --- original text, flown before the rig twin; kept to refill this place:
% \paragraph{RQ2: trajectory-based against dissipative.} On this system the trajectory-based planner was
% better in every test: tracking, detach and attach. The dissipative network did as well as the
% literature suggests only once its tracker was replaced with the velocity loop it assumes, and even then
% it was not robust to a heavier payload. Its strengths are real but did not apply here. With rigid rods
% the payload is held by geometry, which the network can only approximate. Its dissipation acts on
% virtual nodes rather than the payload. And it cannot use the measured payload pose without
% destabilising itself. A combination remains possible, with the planner for tracking and
% reconfiguration and a compliant layer at the weld, but with the OCP resize the evidence does not call
% for it.

\paragraph{Geometry decides what control can do.}
\redo{the geometry cases behind this paragraph: three UAVs at 3, 12 and 9 o'clock hanging the ring tilted, the even ring of four after a detach, the ring on two rim points, and the uneven layout's load share turning into a roll (Sections~\ref{sec:sim_detach} and~\ref{sec:sim_attach}).}
% --- original text, flown before the rig twin; kept to refill this place:
% \paragraph{Geometry decides what control can do.} Several results that first looked like control
% problems were geometry. Three UAVs at 3, 12 and 9 o'clock hang the ring tilted, because its centre of
% mass lies on the edge of their support triangle. An even ring of four loses its support after one
% UAV leaves. A ring on two rim points is a pendulum. An uneven layout puts 40\,\% of the load on one
% UAV, which a tracker without an integrator turns into a roll. Choosing the attachment layout for the
% reconfigurations that will happen is part of the design, and a simple test (is the centre of mass
% inside the support polygon, and by how much) predicts which layouts will work.

\paragraph{Steady errors in a re-planned cascade.} The planner re-solves from the measured state at
every tick, which closes a loop on the payload's position with no integrator. In the twin its
stiffness is about 5\,N/m horizontally and 7\,N/m vertically, so a force of half a newton moves the
payload by 7--10\,cm and stays. The investigation of the payload's height and horizontal offsets
separated four kinds of cause (Table~\ref{tab:error_sources}):
\begin{itemize}
    \item a bias in the cascade itself: the tracker's relaxed initial condition, which held the payload
    about 5\,cm high and which no force estimate could find, because it is not a force;
    \item a wrong belief passed between levels: the breakaway cap, which told the trackers the rods
    pull up to 27\,\% less than they do, and left the rig's ring 20--30\,cm low;
    \item model errors: the rods' mass, one UAV's thrust map;
    \item external forces: a push, wind, and the rig's steady pull towards $-x$.
\end{itemize}
Integral action removes the steady error from all of them, as the internal model principle says, but
only where its gate is open and within its bound. Whether the integral moves the target or acts as a
force in the planner's model changes the transient and what the plan predicts, not the steady state.
A force estimated from motion and thrust removes external forces quickly, but by Newton's law it is
blind to forces the system exerts on itself. A wrong thrust map on one UAV is a moment on the ring,
and only a correction of that UAV's own force removes it. In \cite{tudelft} the onboard incremental
inversion makes each UAV exact at its own force, at 300\,Hz from rotor speed, which is why the planner
needs no integrator there. With Betaflight and a ground-side controller the equivalent was the
measured-cable-force tracker. It levelled the ring against a wrong thrust map, but in this
configuration it upset the landings, so integral action in the planner is the practical substitute.

\paragraph{Simulation and the rig.} The rig twin, flown on exactly the rig's code, matched the rig in
the takeoff (hand-over angle, rod lengths) and in solver health. It also reproduced the rig's height
shortfall from the one setting that caused it. It did not reproduce three things: the slower
breakaway (73--93\,\% of the planned pull against about 97\,\%), the larger tilt near the floor, and
the steady pull towards $-x$. The first is suspected to be an effect near the floor, the second the
downwash, the third the room or its calibration. None of them is modelled. Many apparent controller
faults in simulation turned out to be simulator artefacts, each of which had to be found and removed
before the twin could be trusted:
\begin{itemize}
    \item rate estimates differenced from unstamped poses turned bunched samples into rate spikes under
    CPU load, so the emulated rate loop now uses the gyro;
    \item a proportional-only emulated rate loop could not hold the moment of a rod pivoting below a
    UAV's centre, and flipped UAVs when control was handed over;
    \item the physics engine ignores damping on ball joints, so a damped rod needs universal joints;
    \item wall-clock watchdogs fire when the simulation runs slower than real time, which it does
    (about 0.3 of real time with the full stack on the laptop).
\end{itemize}
Finding them took as long as the controller work.

\paragraph{Limitations.} The work is at lab scale (a 0.86\,kg payload, two to four UAVs, a
6\,$\times$\,6\,m room) with motion capture as the only sensing and all computation on one ground
computer. The rods are rigid, so slack cables are outside the model. The attach depends on an
electromagnet and on a collaborator's approach controller. The hardware results so far cover carrying
only, with three UAVs on 1 October.
% CHECK: stale since the 7-8 Oct rig flights (four UAVs; the detach flown on the rig).
Several simulation tables rest on single runs or on earlier code,
and are marked as such; they will be re-flown on the final code before the results are frozen.
